\documentclass{article}
\usepackage[UTF8]{ctex}
\usepackage{amsmath,amssymb,amsfonts}
\usepackage{enumitem}
\usepackage{geometry}
\geometry{a4paper, margin=1in}
% === 错题本排版优化（不参与题目标识） ===
\setlength{\parindent}{0pt}
\setlength{\parskip}{0.35em plus 0.10em minus 0.08em}
\setlength{\emergencystretch}{3em}
\setlist{itemsep=0.25em, topsep=0.40em, parsep=0pt}
\allowdisplaybreaks[2]
% === 错题本排版优化结束 ===
\begin{document}
\section*{北京工业大学2025年高等代数真题}
\section*{7．北京工业大学2025年高等代数试卷}

\subsection{题目1}
1．计算 \(n\) 阶行列式

\[
\left|\begin{array}{cccccc}
a^{2}+b^{2} & a^{2} b^{2} & 0 & \cdots & 0 & 0 \\
1 & a^{2}+b^{2} & a^{2} b^{2} & \cdots & 0 & 0 \\
0 & 1 & a^{2}+b^{2} & \cdots & 0 & 0 \\
\vdots & \vdots & \vdots & & \vdots & \vdots \\
0 & 0 & 0 & \cdots & a^{2}+b^{2} & a^{2} b^{2} \\
0 & 0 & 0 & \cdots & 1 & a^{2}+b^{2}
\end{array}\right|
\]

\subsection{题目2}
2．设

\[
A=\left(\begin{array}{ccc}
3 & 2 & -2 \\
-k & 1 & k \\
4 & 3 & -3
\end{array}\right) \text { (数据可能有误, 可对比2019 年真题). }
\]

\(k\) 取何值时，存在可逆矩阵 \(P\) ，使得 \(P^{-1} A P\) 为对角矩阵，求 \(P\) 及相应的对角矩阵 \(P^{-1} A P\) 。

\subsection{题目3}
3．证明：欧氏空间中的向量组 \(\alpha_{1}, \alpha_{2}, \cdots, \alpha_{n}\) 线性无关的充要条件是

\[
\left|\begin{array}{cccc}
\left(\alpha_{1}, \alpha_{1}\right) & \left(\alpha_{1}, \alpha_{2}\right) & \cdots & \left(\alpha_{1}, \alpha_{n}\right) \\
\left(\alpha_{2}, \alpha_{1}\right) & \left(\alpha_{2}, \alpha_{2}\right) & \cdots & \left(\alpha_{2}, \alpha_{n}\right) \\
\vdots & \vdots & & \vdots \\
\left(\alpha_{n}, \alpha_{1}\right) & \left(\alpha_{n}, \alpha_{2}\right) & \cdots & \left(\alpha_{n}, \alpha_{n}\right)
\end{array}\right| \neq 0 \text {. }
\]

\subsection{题目4}
4．\(n\) 阶实对称矩阵 \(A\) 的 \(n\) 个特征值为 \(\lambda_{1}, \lambda_{2}, \cdots, \lambda_{n}\) ，且 \(\lambda_{1} \geq \lambda_{2} \geq \cdots \geq \lambda_{n}\) 。证明：

\[
\lambda_{k}=\min _{\operatorname{dim} S=n-k+1} \max _{X \in S, X \neq 0} \frac{(A X, X)}{(X, X)}, k=1,2, \cdots, n .
\]

其中 \(S\) 为向量空间，\(X=\left(x_{1}, x_{2}, \cdots, x_{n}\right)^{\mathrm{T}}\) ，内积（），

\[
(X, Y)=\sum_{i=1}^{n} x_{i} y_{i}
\]

其中 \(Y=\left(y_{1}, y_{2}, \cdots, y_{n}\right)^{\mathrm{T}}\) 。

\subsection{题目5}
5．解答如下问题：

\begin{itemize}
  \item[（1）] \(A\) 是 \(m\) 阶正定矩阵，\(B\) 是 \(m \times n\) 实矩阵，则 \(B^{\mathrm{T}} A B\) 是正定矩阵的充要条件是 \(r(B)=n\) 。
  \item[（2）] \(A\) 是 \(n\) 阶实对称矩阵，证明：\(A\) 为半正定矩阵的充要条件是存在 \(n\) 阶实矩阵 \(B\) ，使得 \(A=B^{\mathrm{T}} B\) 。
\end{itemize}

\subsection{题目6}
6．设 \(A, B\) 为 \(n\) 阶方阵，且 \(A B=A+2025 B\) 。证明：

\begin{itemize}
  \item[（1）] \(r(A)+r(B) \leq r(A B)+n\) ．
  \item[（2）] \(r(A)+r(B) \geq r(A+B)+r(A B)\) ．
\end{itemize}

\subsection{题目7}
7．设 \(\sigma, \tau\) 为 \(n\) 维欧氏空间 \(V\) 中的线性变换，若对任意的 \(\alpha, \beta \in V\) ，都有

\[
(\sigma(\alpha), \beta)=(\alpha, \tau(\beta)) .
\]

则称 \(\sigma\) 是 \(\tau\) 的共轭。证明：

\begin{itemize}
  \item[（1）] \(\sigma\) 是 \(\tau\) 的共轭的充分必要条件是两者在同一组标准正交基下的矩阵互为转置。
  \item[（2）] 若 \(\sigma\) 是 \(\tau\) 的共轭，则 \(\sigma\) 的值域为 \(\tau\) 核空间的正交补，即 \(\operatorname{Im} \sigma=(\operatorname{Ker} \sigma)^{\perp}\) 。
\end{itemize}
\end{document}
